\chapter{Analyse et conception}
\section{Vue d'ensemble de l'architecture}
La solution est organisée comme une chaîne de traitement dont le point d'entrée est la fonction \texttt{run\_pipeline()}. Le choix d'une architecture modulaire répond à deux exigences : adapter les stratégies au type de document et conserver une trace des décisions. Un contrat commun, \texttt{ParsedDocument}, est partagé par les analyseurs ; il contient le texte complet, les métadonnées, les sections structurées et les informations de provenance.

\begin{figure}[H]
\centering
\begin{tikzpicture}[node distance=4mm and 5mm, every node/.style={draw,rounded corners,align=center,minimum height=9mm,font=\small}, >=Latex]
\node[fill=softblue] (doc) {Document\\PDF, TXT, MD};
\node[right=of doc,fill=softblue] (parse) {Ingestion et\\analyse};
\node[right=of parse,fill=softblue] (struct) {Structuration et\\découpage};
\node[right=of struct,fill=softblue] (index) {Indexation et\\recherche};
\node[below=of index,fill=softblue] (extract) {Extraction\\spécialisée};
\node[left=of extract,fill=softblue] (rules) {Normalisation\\et validation};
\node[left=of rules,fill=softblue] (out) {Exports\\JSON / CSV / XLS};
\draw[->,thick] (doc)--(parse); \draw[->,thick] (parse)--(struct); \draw[->,thick] (struct)--(index);
\draw[->,thick] (index)--(extract); \draw[->,thick] (extract)--(rules); \draw[->,thick] (rules)--(out);
\end{tikzpicture}
\caption{Chaîne fonctionnelle de la solution développée}
\label{fig:pipeline}
\end{figure}

La séparation des composants évite qu'une correction de parsing modifie directement la logique d'export. Elle permet également de tester un service de manière isolée, par exemple le découpage des bilans biologiques ou la normalisation des tailles de lésions.

\section{Acteurs et cas d'utilisation}
Les utilisateurs principaux sont l'opérateur qui lance le traitement et le responsable d'étude qui définit les champs attendus. Le système reçoit le document, route son contenu, prépare les candidats et retourne des mises à jour proposées. Un évaluateur hors ligne compare ensuite les sorties avec des gabarits de référence. Cette description représente une assistance au recueil : la décision de validation clinique reste hors du périmètre de la pipeline.

\begin{figure}[H]
\centering
\begin{tikzpicture}[every node/.style={font=\small}, actor/.style={draw,ellipse,fill=softblue,minimum width=2.2cm,minimum height=0.8cm}, use/.style={draw,rounded corners,minimum width=3.6cm,minimum height=0.72cm,align=center}]
\node[actor] (op) at (0,0) {Opérateur};
\node[use] (u1) at (5,1.4) {Traiter un document};
\node[use] (u2) at (5,0) {Consulter les propositions et preuves};
\node[use] (u3) at (5,-1.4) {Exporter les mises à jour};
\node[actor] (study) at (10,0) {Responsable d'étude};
\node[use] (u4) at (5,-2.8) {Configurer le schéma d'étude};
\draw (op)--(u1) (op)--(u2) (op)--(u3); \draw (study)--(u4); \draw[dashed] (u4)--(u1);
\end{tikzpicture}
\caption{Cas d'utilisation principaux}
\end{figure}

\section{Modèle de données et traçabilité}
Le modèle de données est composé d'objets Pydantic. \texttt{ParsedDocument} représente le document analysé. \texttt{DocumentChunk} délimite un extrait et conserve ses positions de caractères. \texttt{ExtractedObservation} porte une observation extraite, tandis que \texttt{FieldCandidate} associe cette observation à une clé canonique et à une valeur normalisée. Enfin, \texttt{EcrfCellUpdate} décrit la mise à jour exportable après validation. Cette progression sépare l'observation brute de la décision de remplissage.

\begin{figure}[H]
\centering
\begin{tikzpicture}[node distance=4mm, every node/.style={draw,rounded corners,align=center,minimum width=3.1cm,minimum height=8mm,font=\small}, >=Latex]
\node[fill=softblue] (a) {ParsedDocument};
\node[below=of a,fill=softblue] (b) {DocumentChunk\\texte + positions};
\node[below=of b,fill=softblue] (c) {ExtractedObservation\\valeur + preuve};
\node[below=of c,fill=softblue] (d) {FieldCandidate\\clé + normalisation};
\node[below=of d,fill=softblue] (e) {EcrfCellUpdate\\champ validé};
\draw[->] (a)--(b); \draw[->] (b)--(c); \draw[->] (c)--(d); \draw[->] (d)--(e);
\end{tikzpicture}
\caption{Progression d'une donnée extraite vers une mise à jour eCRF}
\end{figure}

L'audit enregistre les étapes majeures et la provenance. Cette exigence est essentielle : une valeur affichée sans contexte ni règle de validation ne doit pas être considérée comme directement exploitable.

\section{Stratégie de recherche et d'extraction}
Le système distingue une recherche en mémoire, adaptée au développement et aux tests, d'un backend Qdrant pour un usage vectoriel. Le backend hybride combine représentations denses, recherche BM25 et fusion RRF, avec un réordonnancement optionnel. Cette approche associe la correspondance sémantique à la précision des termes médicaux et des identifiants \cite{qdrant-hybrid}.

Les requêtes ne sont pas globales : elles sont générées à partir du \texttt{StudySchema} actif et ne ciblent que les familles de champs utiles au type documentaire routé. Le schéma déclare notamment la clé canonique, la stratégie d'extraction, la temporalité, la règle de normalisation, la colonne cible et un seuil de confiance.

\section{Contraintes de conception}
La conception a été guidée par plusieurs contraintes. Premièrement, le système doit accepter que toutes les dépendances ne soient pas présentes dans tous les environnements. C'est pourquoi les composants lourds ou optionnels, tels que Docling et les modèles de recherche, ne constituent pas un prérequis pour le chemin de base. Deuxièmement, les opérations d'écriture doivent être réversibles ou reproductibles : les exports peuvent être produits dans une copie du classeur, et la politique de remplacement est explicite.

Troisièmement, l'indexation ne doit pas mélanger les contextes d'études. Le schéma de collections par locataire et les filtres obligatoires sont une réponse technique à ce besoin. Enfin, l'architecture admet que certaines valeurs sont plus sûres lorsqu'elles sont extraites par une règle déterministe, alors que d'autres nécessitent un raisonnement sur le contexte. La pipeline ne force pas une même méthode pour les deux cas.

\section{Diagramme d'activité : décision d'export}
L'activité suivante montre le point où la solution cesse d'être une simple extraction et applique les contrôles nécessaires avant de produire une mise à jour. Un candidat non conforme reste traçable, mais n'est pas inclus dans l'export validé.

\begin{figure}[H]
\centering
\begin{tikzpicture}[node distance=5mm, every node/.style={font=\small,align=center}, process/.style={draw,rounded corners,fill=softblue,minimum width=4.2cm,minimum height=8mm}, decision/.style={draw,diamond,aspect=2,fill=yellow!20,inner sep=1pt}]
\node[process] (a) {Observation extraite};
\node[process,below=of a] (b) {Mapping vers la clé canonique};
\node[process,below=of b] (c) {Normalisation valeur, unité, temporalité};
\node[decision,below=of c] (d) {Conforme au schéma ?};
\node[process,below left=8mm and 14mm of d] (e) {Conserver pour audit\\sans exporter};
\node[decision,below right=8mm and 14mm of d] (f) {Politique d'écriture\\autorisée ?};
\node[process,below=of f] (g) {Créer EcrfCellUpdate};
\draw[->] (a)--(b); \draw[->] (b)--(c); \draw[->] (c)--(d);
\draw[->] (d)--node[left]{non}(e); \draw[->] (d)--node[above right]{oui}(f);
\draw[->] (f)--node[left]{oui}(g); \draw[->] (f.west) |- node[near start,above]{non}(e.east);
\end{tikzpicture}
\caption{Activité de validation précédant l'export}
\end{figure}

\section{Scénario de séquence : extraction d'un bilan}
La figure suivante résume le scénario nominal. Une fois le document fourni, le système l'analyse, reconstruit les lignes biologiques, puis produit des candidats. Les règles métier vérifient les unités et les contraintes du champ avant de générer l'export.

\begin{figure}[H]
\centering
\begin{tikzpicture}[x=2.25cm,y=0.62cm,font=\scriptsize,>=Latex]
\foreach \x/\n in {0/Opérateur,1/Pipeline,2/Parser,3/Extracteur,4/Validation,5/Export} {\node at (\x,0) {\n}; \draw[dashed] (\x,-0.25)--(\x,-7.2);}
\draw[->] (0,-0.8)--(1,-0.8) node[midway,above]{document};
\draw[->] (1,-1.7)--(2,-1.7) node[midway,above]{analyser};
\draw[->] (2,-2.6)--(1,-2.6) node[midway,above]{document structuré};
\draw[->] (1,-3.5)--(3,-3.5) node[midway,above]{chunks ciblés};
\draw[->] (3,-4.4)--(4,-4.4) node[midway,above]{candidats + preuves};
\draw[->] (4,-5.3)--(5,-5.3) node[midway,above]{mises à jour validées};
\draw[->] (5,-6.2)--(0,-6.2) node[midway,above]{JSON / CSV / XLS};
\end{tikzpicture}
\caption{Séquence simplifiée de traitement d'un bilan biologique}
\end{figure}

\section{Conclusion}
La conception repose sur une chaîne contrôlable et configurable. Elle transforme progressivement un document source en proposition d'écriture, sans perdre le lien avec l'évidence. Le chapitre suivant détaille son implémentation.
